% Generated from tex/chapters/ch04/07_実務上の考慮.tex for the SWEBOK structure tree
\subsubsection{構築設計}

構築段階でも設計は行われる。要求分析や上位設計で決まっていない詳細を、構築中に決める必要があるためである。たとえば、アルゴリズム、データ構造、関数やクラスの責務、インタフェース、例外の扱い、ログの粒度などは、構築中に具体化されることが多い。

建築現場で、図面に書かれていない細部を現場で調整する必要があるのと同じで、ソフトウェア構築でも詳細の調整が必要である。ただし、場当たり的に決めるのではなく、設計原則、標準、テスト容易性、変更容易性に基づいて判断する必要がある。
